% Part: computability
% Chapter: computability-theory
% Section: russells-paradox

\documentclass[../../../include/open-logic-section]{subfiles}

\begin{document}

\olfileid{cmp}{thy}{rus}
\olsection{रसेलच्या विरोधापत्तीशी तुलना}

या विभागातील युक्तिवादांची रसेलच्या विरोधापत्तीशी तुलना केल्यास
त्यांतील साम्य आणि भेद स्पष्ट होतात:
\begin{enumerate}
\item रसेलची विरोधापत्ती: $S = \Setabs{x}{x \notin x}$ असे मानू.
  मग $S
  \in S$ असेल तेव्हाच $S \notin S$ असेल; हा विरोध आहे.

  \emph{निष्कर्ष:} असा संच~$S$ अस्तित्वात नाही. ``सर्व संचांचा संच''
  अस्तित्वात आहे हे गृहितक संचसिद्धांताच्या इतर स्वयंसिद्धांशी विसंगत आहे.
%READERNOTE{OLCMP-016 — गोठवलेल्या इंग्रजीत S ∈ S नंतर X ∉ S असे आले आहे, पण X परिभाषित नाही. S च्या व्याख्येवरून आवश्यक विधान S ∉ S आहे; मराठीत तोच निर्देशांक ठेवला आहे.}

\item रसेलच्या विरोधापत्तीचे एक रूपांतर: सर्व फलनांच्या संचापासून
  $\{ 0, 1 \}$ मध्ये जाणारे ``फलन'' $F$ पुढीलप्रमाणे परिभाषित करू:
  \[
  F(f) =
  \begin{cases}
    1 & \text{जर $f$ हे $f$ च्या परिभाषाक्षेत्रात असेल आणि $f(f) = 0$ असेल तर} \\
    0 & \text{अन्यथा}
  \end{cases}
  \]
  तशाच युक्तिवादावरून $F(F) = 0$ असेल तेव्हाच $F(F) = 1$
  असेल; हा विरोध आहे.

  \emph{निष्कर्ष:} $F$ हे फलन नाही. ``सर्व फलनांचा संच'' इतका
  मोठा आहे की तो फलनाचे परिभाषाक्षेत्र असू शकत नाही.

\item विकर्णीकरणाचा युक्तिवाद: $f_0$, $f_1$, \dots हे आंशिक
  संगणनक्षम फलनांचे प्रगणन असू द्या आणि $G \colon \Nat \to
  \{ 0, 1 \}$ ची व्याख्या पुढीलप्रमाणे करू:
  \[
  G(x) =
  \begin{cases}
    1 & \text{जर $f_x(x)\downarrow = 0$ असेल तर} \\
    0 & \text{अन्यथा}
  \end{cases}
  \]
  $G$ संगणनक्षम असेल, तर एखाद्या $k$ साठी ते फलन $f_k$
  असेल. पण मग $G(k) = 1$ असेल तेव्हाच $G(k) = 0$ असेल;
  हा विरोध आहे.

  \emph{निष्कर्ष:} $G$ संगणनक्षम नाही. संचसिद्धांताच्या
  स्वयंसिद्धांनुसार $G$ हे तरीही फलन आहे, हे लक्षात घ्या;
  येथे विरोधापत्ती नाही, तर फरक स्पष्ट झाला आहे.
\end{enumerate}

आंशिक फलने, संगणनक्षम फलने, आंशिक संगणनक्षम फलने इत्यादींबद्दलची
चर्चा गोंधळात टाकू शकते. $\Nat$ पासून $\Nat$ मध्ये जाणाऱ्या सर्व
आंशिक फलनांचा संच हा वस्तूंचा मोठा संग्रह आहे. त्यांपैकी काही
सर्वत्र परिभाषित आहेत, काही आंशिक संगणनक्षम आहेत, काही सर्वत्र
परिभाषित आणि संगणनक्षम अशी दोन्ही आहेत, तर काही यांपैकी कोणतीच नाहीत.
आपण नुसते ``फलन'' म्हटले, तर सामान्यतः सर्वत्र परिभाषित फलन अभिप्रेत
असते, हे लक्षात ठेवा. म्हणून पुढील प्रकार मिळतात:
\begin{enumerate}
\item संगणनक्षम फलने
\item सर्वत्र परिभाषित नसलेली आंशिक संगणनक्षम फलने
\item संगणनक्षम नसलेली फलने
\item सर्वत्र परिभाषितही नसलेली आणि संगणनक्षमही नसलेली आंशिक फलने
\end{enumerate}
हा फरक समजण्यासाठी $\Nat$ पासून $\Nat$ मध्ये जाणारी सर्व आंशिक
फलने दर्शवणारा मोठा चौकोन काढा. त्यात सर्वत्र परिभाषित फलने
आणि आंशिक संगणनक्षम फलने दर्शवणारी, एकमेकांवर काही प्रमाणात
येणारी दोन क्षेत्रे दाखवा. अशा आकृतीत मिळणाऱ्या प्रत्येक
क्षेत्रातील एखाद्या वस्तूचे वर्णन करता येते का, हे पाहणे चांगला सराव आहे.

\end{document}
